\section{Conclusion}
This Systematic Literature Review set out to answer the following research question: \textit{"What are the dominant paradigms, open challenges, and emerging research directions in cyber deception, particularly in the shift from static/reactive honeypots to adaptive, AI-driven and predictive deception systems?"}

Our analysis of the literature confirms that the field of cyber deception has undergone, and continues to undergo, a profound evolution. The dominant paradigm has decisively shifted from passive detection to proactive, intelligent adversary engagement.

The evolutionary path, as mapped in our analysis, is clear. The discipline's foundations were built on \textbf{static, reactive traps}---primarily honeypots---whose main purpose was detection, forensic analysis, and intelligence gathering \cite{spitzner2002honeypots, provos2004virtual, nawrocki2016survey}. This model was subsequently formalized and expanded by comprehensive \textbf{taxonomies} \cite{han2018deception} and sophisticated \textbf{proactive strategies}, such as those informed by game theory \cite{pawlick2019game, carroll2011game} and psychological principles \cite{ferguson2021examining}.

Today, the field is defined by an escalating "arms race" driven by Artificial Intelligence. The most dominant emerging paradigm is that of \textbf{AI-driven, adaptive deception}. These systems, often orchestrated by Reinforcement Learning (RL), are no longer static; they are predictive, dynamic, and capable of tailoring deception policies in real-time \cite{hammad2024deep}. This evolution is a direct and necessary response to the emergence of \textbf{AI-powered adversaries}, which leverage LLMs to automate and scale sophisticated attacks \cite{deng2023pentestgpt, fang2024llm}.

Our review concludes that the most significant \textbf{open challenges} exist at this new frontier. Future research must focus not just on building technically proficient decoys, but on solving the meta-problems of scalability, developing new metrics for proactive success, and, most critically, designing deceptions capable of fooling autonomous AI agents. The era of passive cyber defense is over. The future of deception lies in creating intelligent, adaptive systems that can "play the game" as actively and intelligently as the adversaries they face.